% Код LaTeX, переписывающий текст "Метод Зейделя" 1 в 1
\documentclass[12pt]{article}
\usepackage[utf8]{inputenc}
\usepackage[T2A]{fontenc}
\usepackage[russian]{babel}
\usepackage{amsmath,amssymb}
\usepackage{geometry}
\geometry{a4paper, margin=2.5cm}

\begin{document}
дополнительное задание цуканников кирилл \\
оригинал: \\
\section*{Метод Зейделя}

Полагая $B = A_{L} + D, \tau = 1$, получим

\begin{equation*}
(A_{L} + D)\bigl(x^{(k+1)} - x^{(k)}\bigr) + (A_{L} + D + A_{V})x^{(k)} = f, \quad k = 0,1,\dots
\end{equation*}

который называется методом Зейделя. Для очередного приближения в методе Зейделя приходится решать систему уравнений

\begin{equation*}
(A_{L} + D)x^{(k+1)} = f - A_{V}x^{(k)}.
\end{equation*}

Так как $A_{L} + D$ - треугольная нижняя матрица, то явные формулы:

\begin{equation*}
x_{i}^{(k+1)} = \frac{1}{a_{ii}}\Bigl( f_{i} - \sum_{j=i+1}^{n} a_{ij} x_{j}^{(k)} - \sum_{j=1}^{\,i-1} a_{ij} x_{j}^{(k+1)} \Bigr), \quad i = 1,\dots,n.
\end{equation*}
\\
КАК ПОЛУЧИЛИ ПОКОМПОНЕНТУЮ ФОРМУЛУ: \\
$Ax = f$ \\
$ A = A_L + D + A_V $ \\
$A_L$ - строго нижнетреуголньая матрица \\
$D$ - диагональная матрица (на диагонали $a_{ii}$)  \\
$A_V$ - строго верхнетреуголньая матрица \\
упрощенно можно записать что $A = B - (B - A)$ , подставляю в Ax = f \\
$( B - (B - A))x = f $  и перехожу к итерациям \\
\[
B x^{(k+1)} = (B-A)\,x^{(k)} + f.
\]
В самом метаде Зейделя пусть $B = A_L + D$ - легко считать методом простой итерации \\
тогда подставляю и получаю $ B - A = A_L + D  - (A_L+D+A_V)  = -A_V$ \\
группирую и получаю: $ (A_L+D)x^{(k+1)} = f-A_Vx^{(k)} $\\
теперь из $ A = A_L + D + A_V $ получаю $A_V = A-(A_L+D)$, подставляю в то что получил выше и получаю \\
$(A_L+D) x^{(k+1)} = f - (A-(A_L+D))x^{(k)} $ \\
$(A_L+D) ( x^{(k+1)} - x^{(k)}) + Ax^{(k)} = f $ \\

\begin{equation*}
 Ax = \sum_{j=1}^{n} a_{ij} x_{j} = f_i; i = 1, ..., n \\
\end{equation*}
$A_L$ - элементы с j<i \\
$A_V$ - элементы с j>i \\
$ (A_L+D)x^{(k+1)} = f-A_Vx^{(k)} $  - записываем в виде суммы и записываем i-ю строку\\

\begin{equation*}
\sum_{j=1}^{i-1} a_{ij} x_{j}^{(k+1)} + a_{ii}x_i^{(k+1)} = f_i - \sum_{j=i+1}^{n} a_{ij}x_j^{(k)}\\
\end{equation*}
$  a_{ii} $ - то что осталось от диагональной матрицы на iй строке \\
теперь можно второе слогаемое вынести и разделить на $a_{ii}$ \\
\begin{equation*}
x_{i}^{(k+1)} = \frac{1}{a_{ii}}\Bigl( f_{i} - \sum_{j=i+1}^{n} a_{ij} x_{j}^{(k)} - \sum_{j=1}^{\,i-1} a_{ij} x_{j}^{(k+1)} \Bigr), \quad i = 1,\dots,n.
\end{equation*}
это и будет покомнонетная формула метода Зейделя.\\
РАЗБОР ОШИБКИ КОПИПАСТОМ \\
в лекции у $\sum_{j=1}^{\,i-1} a_{ij} x_{j}^{(k+1)}$ было записано не (k+1) а просто (k). хотя по вычислениям этого не может быть и по методу Зейделя тоже. на момент i мы уже посчитали включительно до $x_{i-1}^{(k+1)}$. 

\end{document}
